\chapter{PID 控制器设计} \label{ch:pid-control}

很多单片机程序都要回答同一个问题：测量值偏离目标以后，执行器应该输出多大的控制量？温度控制、电机调速与云台角度控制看起来各不相同，但它们都可以抽象成“测量偏差，再根据偏差修正输出”的闭环系统。本章从反馈控制的基本结构入手，依次介绍比例、积分与微分三种作用，再实现带积分限幅、输出限幅和微分滤波的 PID 控制器。完整示例代码位于 \texttt{Chapter-pid-control/PID控制/} 目录。

\medskip 本章涉及到的实验内容：
\begin{enumerate}
\item 闭环控制系统与误差的定义；
\item 位置式 PID 与增量式 PID；
\item 积分饱和、输出限幅与微分噪声处理；
\item 一阶惯性对象的闭环控制实验。
\end{enumerate}

% ---------------------------------------------------------------------
\section{从开环控制到闭环控制}

开环控制只负责给出命令，不关心被控对象有没有达到目标。例如固定输出一半占空比驱动电机，电池电压、机械负载和摩擦力发生变化以后，转速也会跟着变化。闭环控制则不断读取传感器反馈，用设定值与反馈值之差修正控制量。

设定值记为 $r(t)$，反馈值记为 $y(t)$，两者的误差为：

\begin{equation}
e(t)=r(t)-y(t)
\label{eq:pid-error}
\end{equation}

图~\ref{fig:pid-control-loop}~给出了一个典型的 PID 闭环控制系统。圆形节点完成误差计算，PID 控制器根据误差产生控制量 $u(t)$，被控对象输出 $y(t)$，传感器再把输出送回比较节点。

\begin{figure}[htbp]
\centering
\begin{tikzpicture}[
    block/.style={draw=structurecolor, thick, rounded corners=2pt, minimum width=2.7cm, minimum height=1.0cm, align=center, fill=structurecolor!8},
    sum/.style={draw=structurecolor, thick, circle, minimum size=0.8cm},
    arrow/.style={-stealth, thick}
]
    \node (setpoint) at (-5.2,0) {$r(t)$};
    \node[sum] (sum) at (-3.7,0) {$+$};
    \node[block] (pid) at (-1.2,0) {PID 控制器};
    \node[block] (plant) at (2.2,0) {被控对象};
    \node (output) at (5.0,0) {$y(t)$};
    \node[block, minimum width=2.3cm] (sensor) at (2.2,-2.0) {反馈测量};

    \draw[arrow] (setpoint) -- (sum);
    \draw[arrow] (sum) -- node[above] {$e(t)$} (pid);
    \draw[arrow] (pid) -- node[above] {$u(t)$} (plant);
    \draw[arrow] (plant) -- (output);
    \draw[arrow] (4.1,0) |- (sensor.east);
    \draw[arrow] (sensor.west) -| node[pos=0.9,left] {$-$} (sum.south);
\end{tikzpicture}
\caption{PID 闭环控制系统框图} \label{fig:pid-control-loop}
\end{figure}

闭环不等于一定稳定。控制器参数过小，系统响应迟缓；参数过大，输出会在目标附近反复振荡。PID 参数整定的本质，就是在响应速度、稳态误差与抗干扰能力之间寻找平衡。

% ---------------------------------------------------------------------
\section{PID 三种控制作用}

\subsection{比例作用}

比例项直接放大当前误差：

\begin{equation}
u_{P}(t)=K_{p}e(t)
\end{equation}

$K_p$ 越大，系统对误差的反应越强，但过大的比例系数也容易引起超调和振荡。只使用比例控制时，为了维持一个非零控制输出，系统往往需要保留一部分稳态误差。

\subsection{积分作用}

积分项累积从开始运行到当前时刻的全部误差：

\begin{equation}
u_{I}(t)=K_{i}\int_{0}^{t}e(\tau)\,\mathrm{d}\tau
\end{equation}

只要误差还没有归零，积分项就会继续变化，因此它能够消除稳态误差。代价是积分存在记忆，执行器已经达到极限时仍继续累积，就会形成积分饱和。

\subsection{微分作用}

微分项观察误差的变化速度：

\begin{equation}
u_{D}(t)=K_{d}\frac{\mathrm{d}e(t)}{\mathrm{d}t}
\end{equation}

误差正在快速增大时，微分项提前加强修正；误差正在快速减小时，微分项提前收住控制量。因此微分具有预测趋势、抑制超调的作用。然而差分运算也会放大测量噪声，工程中通常需要对微分项进行低通滤波。

综合三种作用，连续形式的 PID 为：

\begin{equation}
u(t)=K_{p}e(t)+K_{i}\int_{0}^{t}e(\tau)\,\mathrm{d}\tau+K_{d}\frac{\mathrm{d}e(t)}{\mathrm{d}t}
\label{eq:continuous-pid}
\end{equation}

% ---------------------------------------------------------------------
\section{离散化与程序实现}

单片机只能在离散时刻计算控制量。设固定采样周期为 $T_s$，第 $k$ 次采样的误差为 $e[k]$，位置式 PID 可以写成：

\begin{align}
I[k] &= I[k-1]+K_iT_se[k] \label{eq:pid-integral}\\
D[k] &= K_d\frac{e[k]-e[k-1]}{T_s} \label{eq:pid-derivative}\\
u[k] &= K_pe[k]+I[k]+D[k] \label{eq:pid-position}
\end{align}

这里特意使用 \texttt{align} 环境排列多条相互关联的公式。正文中的单个变量使用行内数学环境，独立且需要编号的公式使用 \texttt{equation} 环境，程序名称和路径则使用 \texttt{texttt} 命令，三种内容不要混在一起。

示例代码用 \texttt{PIDController\_T} 保存参数和历史状态。公共函数包括：

\begin{itemize}
\item \texttt{PID\_Init}：设置控制周期与工作模式；
\item \texttt{PID\_SetTunings}：设置 $K_p$、$K_i$ 和 $K_d$；
\item \texttt{PID\_SetLimits}：设置积分项和最终输出的范围；
\item \texttt{PID\_SetDerivativeFilter}：设置微分低通滤波系数；
\item \texttt{PID\_Calculate}：完成一个控制周期的计算。
\end{itemize}

调用方式如下：

\begin{lstlisting}[style=c]
PIDController_T controller;

PID_Init(&controller, PID_MODE_POSITIONAL, 0.01f);
PID_SetTunings(&controller, 1.8f, 0.8f, 0.05f);
PID_SetLimits(&controller, -30.0f, 30.0f, -100.0f, 100.0f);
PID_SetDerivativeFilter(&controller, 0.2f);

/* 固定周期读取反馈，并计算新的执行器输出 */
controlOutput = PID_Calculate(&controller, setpoint, feedback);
\end{lstlisting}

\begin{anybox}{固定控制周期}
PID 中的积分与微分都依赖采样周期。如果 \texttt{PID\_Calculate} 的调用间隔忽长忽短，等效参数也会跟着变化。裸机程序应把控制计算放进固定周期定时任务，而不是放在执行时间不确定的主循环里。
\end{anybox}

% ---------------------------------------------------------------------
\section{位置式与增量式 PID}

位置式 PID 直接给出当前控制量，结构直观，适合温度和位置控制。增量式 PID 则计算本周期输出相对于上一周期的变化量：

\begin{equation}
\begin{aligned}
\Delta u[k]={}&K_p\bigl(e[k]-e[k-1]\bigr)+K_iT_se[k]\\
&+\frac{K_d}{T_s}\bigl(e[k]-2e[k-1]+e[k-2]\bigr)
\end{aligned}
\label{eq:pid-increment}
\end{equation}

再把增量累加到上一周期输出：

\begin{equation}
u[k]=u[k-1]+\Delta u[k]
\end{equation}

增量式 PID 在切换控制器时冲击较小，常用于电机速度控制。它依赖上一周期输出，因此执行器输出被其他逻辑修改以后，需要同步更新控制器状态，或者调用 \texttt{PID\_Reset} 重新开始。

% ---------------------------------------------------------------------
\section{工程中的三个保护措施}

\subsection{积分限幅}

执行器达到最大输出以后，继续累积积分并不能让执行器更用力，只会使控制器恢复得更慢。示例把积分状态限制在 $I_{\min}$ 和 $I_{\max}$ 之间：

\begin{equation}
I[k]=\operatorname{clip}\bigl(I[k],I_{\min},I_{\max}\bigr)
\end{equation}

\subsection{输出限幅}

计算结果还要受到执行器物理能力约束。例如占空比只能在允许范围内变化，双向电机控制量则可以设置为对称的正负区间。限幅应放在控制器内部，避免任何调用者忘记处理。

\subsection{微分滤波}

示例先计算原始误差微分 $d_{\mathrm{raw}}[k]$，再用一阶低通滤波得到 $d_f[k]$：

\begin{equation}
d_f[k]=d_f[k-1]+\alpha\bigl(d_{\mathrm{raw}}[k]-d_f[k-1]\bigr)
\end{equation}

$\alpha$ 越小，微分输出越平滑，但相位延迟也会增加。噪声较大的编码器和 ADC 信号可以从较小系数开始调整。

% ---------------------------------------------------------------------
\section{控制响应示意图}

图~\ref{fig:pid-step-response}~展示了同一个阶跃目标下三种典型响应。该图用于理解整定方向，并不是某一组参数的精确仿真结果。

\begin{figure}[htbp]
\centering
\begin{tikzpicture}[x=1.05cm,y=0.8cm]
    \draw[->,thick] (0,0) -- (10.5,0) node[right] {时间};
    \draw[->,thick] (0,0) -- (0,6.2) node[above] {输出};
    \draw[dashed,thick] (0,4.5) -- (10,4.5) node[right] {设定值};
    \draw[structurecolor,thick] plot[domain=0:10,samples=100] (\x,{4.5*(1-exp(-0.42*\x))});
    \draw[red,thick] plot[domain=0:10,samples=120] (\x,{4.5*(1-exp(-0.55*\x)*(cos(deg(1.5*\x))+0.25*sin(deg(1.5*\x))))});
    \draw[green!50!black,thick] plot[domain=0:10,samples=100] (\x,{4.2*(1-exp(-0.8*\x))});
    \node[structurecolor] at (7.8,2.8) {响应迟缓};
    \node[red] at (6.8,5.4) {超调振荡};
    \node[green!50!black] at (8.0,3.9) {存在稳态误差};
\end{tikzpicture}
\caption{PID 参数不同时的阶跃响应示意图} \label{fig:pid-step-response}
\end{figure}

参数整定可以从纯比例控制开始：先把 $K_i$ 和 $K_d$ 设为零，逐步增加 $K_p$，直到响应速度合适但还没有持续振荡；随后增加 $K_i$ 消除稳态误差；最后少量增加 $K_d$ 抑制超调。每次只调整一个参数，才能看清变化来自哪里。

% ---------------------------------------------------------------------
\section{PC 验证结果}

本章示例使用 C99 编译，并打开 \texttt{-Wall}、\texttt{-Wextra}、\texttt{-Wpedantic} 和 \texttt{-Werror} 严格警告选项。\texttt{pid\_control.c} 与 \texttt{main.c} 编译过程没有产生警告，程序正常结束。

示例以 $10\,\mathrm{ms}$ 为控制周期，共运行 200 次，因此表~\ref{tab:pid-pc-result}~中的末次结果对应 $2\,\mathrm{s}$。设定值为 $50$，被控对象采用一阶惯性模型。

\begin{table}[htbp]
\centering
\caption{PID 示例的 PC 运行结果} \label{tab:pid-pc-result}
\begin{tabular}{llll}
\toprule
迭代序号 & 设定值 & 对象输出 & 控制输出 \\
\midrule
199 & 50.000 & 41.492 & 44.819 \\
\bottomrule
\end{tabular}
\end{table}

两秒仿真结束时，对象输出仍在向设定值上升，因此该结果不能解释为稳态值。这个实验主要验证控制器能够连续更新、控制输出没有越过设定限幅，并且对象输出沿正确方向变化。需要观察稳态误差时，应延长运行时间，或者根据被控对象重新整定参数。

除示例程序外，还使用独立断言验证了纯比例位置式 PID、纯比例增量式 PID、非法控制周期和空指针处理，断言均正常通过。

% ---------------------------------------------------------------------
\section{实验报告：PID 控制器移植与整定}

\subsection{实验基本信息}

MCU 型号：\LabFillLine[3.5cm] \quad 工作主频：\LabFillLine[3.5cm]

控制周期：\LabFillLine[3.5cm] \quad 被控对象类型：\LabFillLine[3.5cm]

比例系数 $K_p$：\LabFillLine[3.5cm] \quad 积分系数 $K_i$：\LabFillLine[3.5cm]

微分系数 $K_d$：\LabFillLine[3.5cm] \quad 输出限幅范围：\LabFillLine[3.5cm]

\subsection{实验目的}
\begin{enumerate}
\item 掌握位置式 PID 与增量式 PID 在单片机上的移植与参数配置方法；
\item 理解比例、积分、微分三种控制作用对系统响应速度、稳态误差和超调量的影响；
\item 学会通过阶跃响应曲线整定 PID 参数，能够根据实测波形判断参数调整方向；
\item 掌握积分限幅、输出限幅和微分滤波等工程保护措施的作用与配置方法。
\end{enumerate}

\subsection{实验原理}
PID 控制器根据设定值与反馈值的偏差 $e[k]$ 计算控制量。位置式 PID 直接输出当前控制量：$u[k] = K_p e[k] + I[k] + D[k]$，其中积分项 $I[k] = I[k-1] + K_i T_s e[k]$，微分项 $D[k] = K_d (e[k]-e[k-1])/T_s$。工程中需对积分项和输出进行限幅，并对微分项进行低通滤波以抑制噪声放大。增量式 PID 则计算输出变化量 $\Delta u[k]$，适合执行器需要增量控制的场景。参数整定通常从比例开始，再逐步引入积分和微分。

\subsection{硬件连接}

\LabPasteBox{PID 控制硬件连接}{粘贴开发板与执行器（如电机驱动或加热器）、传感器（编码器、ADC）以及调试串口的连接图片，标注电源、PWM 输出、反馈信号和公共地。}{fig:ch13-pid-hardware}

\subsection{关键代码}

\LabPasteBox{PID 控制器初始化、更新与参数整定代码}{粘贴 PID 结构体定义、参数设置函数、控制周期定时任务和 \texttt{PID\_Calculate} 调用代码，并标注积分限幅和输出限幅的配置。}{fig:ch13-pid-code}

\subsection{实验步骤}
\begin{enumerate}
\item 将 `pid\_control.c` 与 `pid\_control.h` 加入单片机工程，确认控制周期定时器已正确配置（如 10ms）。
\item 连接传感器和执行器，编写反馈读取函数和执行器输出函数，将原始 ADC 值或编码器计数值换算为具有实际物理单位的反馈值。
\item 从纯比例控制开始：将 $K_i$ 和 $K_d$ 设为 0，逐步增大 $K_p$，观察阶跃响应，直到响应速度合适且未出现持续振荡。
\item 加入积分项：缓慢增大 $K_i$，观察稳态误差是否逐渐消除，注意积分限幅是否触发。
\item 加入微分项：少量增加 $K_d$，观察超调量是否减小，若反馈噪声较大则启用微分滤波并调整滤波系数。
\item 通过串口以固定周期输出设定值、反馈值和控制量，使用上位机绘图工具（如 SerialPlot、Python）绘制响应曲线。
\item 依次测试三种场景：① 阶跃响应（设定值突变）；② 负载扰动（施加外部干扰）；③ 设定值切换（多级目标变化），记录各场景下的响应指标。
\item 测试输出限幅和积分限幅的保护效果：将输出限幅设为一个较小值，观察控制量是否被正确截断且积分项未过度累积。
\end{enumerate}

\subsection{实验现象与数据记录}

\begin{table}[htbp]
\centering
\caption{PID 控制器阶跃响应与扰动实验记录}
\label{tab:ch13-pid-record}
\begin{tabularx}{\textwidth}{p{2.4cm}p{2.2cm}p{2.0cm}p{2.0cm}p{2.0cm}X}
\toprule
测试场景 & $K_p$ / $K_i$ / $K_d$ & 上升时间(ms) & 超调量(\%) & 稳态误差 & 备注 \\
\midrule
纯比例 & \LabFillLine[1.8cm] & \LabFillLine[1.6cm] & \LabFillLine[1.6cm] & \LabFillLine[1.6cm] & \LabFillLine[1.6cm] \\
比例+积分 & \LabFillLine[1.8cm] & \LabFillLine[1.6cm] & \LabFillLine[1.6cm] & \LabFillLine[1.6cm] & \LabFillLine[1.6cm] \\
比例+积分+微分 & \LabFillLine[1.8cm] & \LabFillLine[1.6cm] & \LabFillLine[1.6cm] & \LabFillLine[1.6cm] & \LabFillLine[1.6cm] \\
负载扰动恢复 & \LabFillLine[1.8cm] & \LabFillLine[1.6cm] & \LabFillLine[1.6cm] & \LabFillLine[1.6cm] & \LabFillLine[1.6cm] \\
\bottomrule
\end{tabularx}
\end{table}

\LabPasteBox{PID 阶跃响应与扰动波形}{粘贴串口绘图工具绘制的设定值、反馈值和控制量曲线截图，标注上升时间、超调量和稳态误差。}{fig:ch13-pid-result}

\subsection{结果分析}

纯比例控制下是否存在稳态误差及其大小：\LabFillLine[6cm]

加入积分项后稳态误差的变化：\LabFillLine[6cm]

加入微分项后超调量和上升时间的变化：\LabFillLine[6cm]

输出限幅触发时，积分项是否继续累积（积分饱和现象）：\LabFillLine[6cm]

\subsection{实验总结}

\noindent\LabFillLine[\linewidth]\par
\noindent\LabFillLine[\linewidth]\par
\noindent\LabFillLine[\linewidth]

% ---------------------------------------------------------------------
\section{本章小结}

PID 并不神秘：比例项处理当前误差，积分项处理历史误差，微分项观察误差趋势。真正决定工程可靠性的往往不是公式本身，而是固定采样周期、积分限幅、输出限幅和噪声处理。把这些边界条件封装进控制器以后，同一套代码就能稳定地复用到不同被控对象中。

% ---------------------------------------------------------------------
\begin{problemset}
\item 比例、积分和微分三种控制作用分别主要改善闭环系统的什么特性？为什么微分项通常不直接对含噪声的原始误差作差分？
\item 某位置式 PID 的参数为 $K_p=2$、$K_i=0.5$、$K_d=0.1$，采样周期为 $T_s=0.1\,\mathrm{s}$。已知上一周期积分状态为 $I[k-1]=1$，当前误差为 $e[k]=3$，上一周期误差为 $e[k-1]=2$。忽略限幅与微分滤波，计算当前积分项、微分项和控制输出。
\item 积分限幅与输出限幅解决的问题有何不同？如果执行器已经达到输出上限，但积分项仍持续累加，系统恢复时可能出现什么现象？
\item 位置式 PID 与增量式 PID 的输出含义有什么区别？当执行器输出被手动控制逻辑改写以后，为什么增量式控制器需要同步内部状态或执行复位？
\item 将本章 PID 控制器移植到 STM32 电机调速系统时，请设计从“定时采样”到“安全输出”的最小数据流程，并说明控制周期、反馈异常和输出限幅应分别在哪里处理。
\end{problemset}
